\section{Algorithmic Composition}\label{sec:algorithmic-composition}

Algorithmic composition uses formal rules to generate music rather than placing every note by hand.
Computer-assisted approaches date to the \textit{MUSIC} programs at Bell Labs~\cite{Roads_1996}; Nierhaus surveys
rule- and process-driven generation more broadly~\cite{Nierhaus_2009}.
MIDI (1983) later standardized symbolic event exchange~\cite{Moog_1986}, while live-coding languages such as
TidalCycles and Sonic Pi emphasized real-time performance over offline artifacts~\cite{McLean_Wiggins_2010,Aaron_2016}.

The practical advantage for this thesis is \textit{source-level compression}: a pattern invoked inside a loop
describes repetition declaratively, and the compiler expands it into a flat event timeline at build time.
Listing~\ref{lst:compression-example} sketches the idea; later material specifies and implements the mechanism in full.

\begin{lstlisting}[label={lst:compression-example},caption={Compact source that lowering expands into many MIDI events.}]
pattern motif(int dur) {
    play A4 :dur;
    play rest :dur;
    play C5 :dur;
}

track melody using piano {
    loop (256) { play motif(1); }
}
\end{lstlisting}
